\documentclass[11pt,a4paper]{article}

\usepackage[utf8]{inputenc}

\usepackage[T1]{fontenc}

\usepackage[polish]{babel}

\usepackage{lmodern}

\usepackage{geometry}

\geometry{a4paper,margin=27mm}

\usepackage{microtype}

\usepackage[hidelinks]{hyperref}

% Portfolio: MSC-AR-P09; source: prospective ARC-XR research scope.

\title{Wykrywanie błądzenia myślami i odzyskiwanie uwagi z EEG i okulografii}

\author{Mateusz Pomianek}

\date{Wersja 0.3.0, 21 listopada 2026}

\begin{document}

\maketitle

\begin{abstract}

Przedmiotem opracowania jest Wykrywanie błądzenia myślami i odzyskiwanie uwagi z EEG i okulografii. Czy połączenie EEG i okulografii pozwala wykrywać epizody błądzenia myślami w zadaniu? Porównane zostaną modele oparte na EEG, spojrzeniu i obu modalnościach, z wykorzystaniem niezależnie pozyskanych znaczników epizodów uwagi. Głównym punktem końcowym jest trafność wykrywania epizodów przy określonym czasie reakcji i udziale próbek odrzuconych. Opis ma charakter wstępny i nie przedstawia niezweryfikowanych wyników.

\end{abstract}

\section{Wprowadzenie}

Rozwój adaptacyjnych systemów XR stwarza możliwość badania interakcji człowieka z informacją przestrzenną, ale samo zastosowanie AR, VR lub modeli sztucznej inteligencji nie przesądza o przewadze poznawczej czy ergonomicznej. W rozpatrywanym zagadnieniu, Wykrywanie błądzenia myślami i odzyskiwanie uwagi z EEG i okulografii, kluczowe jest rozdzielenie właściwości systemu od zmienności zachowania użytkownika. Porównanie należy zaprojektować tak, aby mierzyć jednoznaczne zjawisko, a nie tylko ogólną ocenę technologii.

\section{Cel pracy i pytanie badawcze}

Czy połączenie EEG i okulografii pozwala wykrywać epizody błądzenia myślami w zadaniu? Projekt wymaga krytycznego przeglądu aktualnego stanu badań oraz zdefiniowania nowości w stosunku do innych tematów portfolio XR.

\section{Zakres i odrębny wkład naukowy}
Predykcja epizodów błądzenia myślami na podstawie sygnałów EEG i spojrzenia przy niezależnych znacznikach uwagi; odzyskiwanie uwagi jako interwencja pozostaje oddzielnym badaniem. Odrębność to trafność czasowego rozpoznania deklarowanych epizodów, nie diagnostyka lub twierdzenie o obiektywnym odczycie myśli.
\section{Hipotezy badawcze}
\paragraph{H1.} Fuzja EEG i wzroku zwiększa AUPRC epizodów błądzenia myślami wobec najsilniejszego modelu jednej modalności.
\paragraph{H2.} Wynik jest dodatni po kontrolowaniu jakości zgłoszeń epizodów i liczby pominiętych próbek.
\paragraph{H3.} Różnica utrzymuje się na sesjach oddzielonych od treningu, ale maleje wraz z przesunięciem profilu osoby.
\section{Projekt badania i plan eksperymentów}
\subsection{A. Operacjonalizacja i kwalifikacja danych}
Predykcja epizodów błądzenia myślami na podstawie sygnałów EEG i spojrzenia przy niezależnych znacznikach uwagi; odzyskiwanie uwagi jako interwencja pozostaje oddzielnym badaniem. Jednostka wnioskowania: osoba × niezależna sesja i epizod uwagi. W protokole zamrozić definicję i mianownik wyniku oraz źródło prawdy odniesienia, zgodność znaczników czasu, prawa do danych i zasady rekrutacji.
\subsection{B. Eksperyment porównawczy}
Pozyskiwać nienaprowadzające próby sondujące uwagę w zadaniu ekranowym z poprawnym synchronizowaniem EEG i Tobii Eye Tracker 5; nie deklarować pomiaru HMD. Estymand pierwszorzędowy: różnica AUPRC między modelem fuzji EEG+spojrzenie a silniejszym jednosygnałowym modelem na niewidzianych sesjach. Przed testem zdefiniować randomizację, próbę, punkt końcowy i warunki wykluczenia bez selekcji wyniku post hoc.
\subsection{C. Kontrole i analiza odporności}
Zmniejszyć ryzyko przecieku okien EEG ze zbioru jednej osoby; ustalić opóźnienia raportu i dopuszczalne okno definicyjne epizodu. Zebrać odrębne sesje, kontrolę losową etykiet, stan zmęczenia i poziom braków obu strumieni. W rejestrze uwzględnić braki, wstrzymania, opóźnienia i niepełne próby zamiast usuwać je z analizy.
\subsection{D. Walidacja na niezależnym materiale}
Weryfikować na niezależnych sesjach lub osobach, bez dostępu modelu do danych końcowych. Porównanie uczestnikowo-klastrowym bootstrapem AUPRC i odsetka odrzuconych próbek; losowy podział klatek zakazany. Dokumentować wyniki niesprzyjające hipotezie oraz rzeczywiste granice przeniesienia.
\section{Estymand, analiza statystyczna i decyzje}
różnica AUPRC między modelem fuzji EEG+spojrzenie a silniejszym jednosygnałowym modelem na niewidzianych sesjach. Jednostka: osoba × niezależna sesja i epizod uwagi. Porównanie uczestnikowo-klastrowym bootstrapem AUPRC i odsetka odrzuconych próbek; losowy podział klatek zakazany. H1 jest kandydatem hipotezy potwierdzającej, a H2/H3 pozostają eksploracyjne do formalnego planu wielokrotnych porównań. Liczebność uczestników, margines, próg i moc lub precyzja wymagają pilotażu oraz zamrożenia przed testem końcowym.
\section{Stan badań i granica nowości}
\begin{itemize}
\item EEG/VR time alignment and experimental validity must be explicitly documented before any multimodal inference claim. \cite{xrprior1}.
\item Eye tracking in immersive VR requires explicit validity and construct assumptions; device integration cannot be presumed. \cite{xrprior2}.
\end{itemize}
Odrębność to trafność czasowego rozpoznania deklarowanych epizodów, nie diagnostyka lub twierdzenie o obiektywnym odczycie myśli. Odrębność od wcześniejszych publikacji i innych tematów programu należy sprawdzić przed rejestracją badania.
\section{Warunki bezpieczeństwa, sprzętu i zgodności}
Jest to plan prospektywny, nie dowód wykonania eksperymentu lub zatwierdzonej etyki. Udział osób, rejestrowanie biometrii i interwencje wymagają uprzedniej procedury etycznej i świadomej zgody. Nie wnioskować o emocjach ani zdrowiu z samych sygnałów. Standardowy Quest 3 nie ma okulometrii HMD, a Tobii Eye Tracker 5 służy do stanowiska ekranowego; każde połączenie zweryfikować przed protokołem.
\section{Materiał i metody}

Porównane zostaną modele oparte na EEG, spojrzeniu i obu modalnościach, z wykorzystaniem niezależnie pozyskanych znaczników epizodów uwagi. Wymagane jest zdefiniowanie jednostki analizy, zasad doboru przypadków, wersji scenariusza i źródeł pomiaru. Rozkład warunków oraz reguły wykluczenia prób należy określić przed pomiarem. W badaniach z udziałem ludzi niezbędna jest uprzednia właściwa ocena etyczna i świadoma zgoda. W testach wyłącznie symulacyjnych należy oddzielić prawdę symulacyjną od obserwacji empirycznej.

\subsection{Wynik główny i kryteria jakości}

Trafność wykrywania epizodów przy określonym czasie reakcji i udziale próbek odrzuconych. Pomiar powinien zachowywać jawne jednostki, znaczniki czasu i reguły oceny poprawności. Przy modelach predykcyjnych stosować rozłączne partycje oceny; przy eksperymentach interwencyjnych oddzielić skuteczność predykcji od przyczynowego wpływu interfejsu.

\section{Plan analizy}

Walidacja między uczestnikami i sesjami, analiza braków sygnału i niepewności znaczników czasu. W wynikach należy raportować wielkość efektu lub błąd estymacji z przedziałem niepewności, liczbą obserwacji i odsetkiem odrzuconych prób. Analizy dodatkowe wymagają jasnego oznaczenia jako eksploracyjne. Nie wybierać progów decyzyjnych ani modelu na podstawie obserwacji ze zbioru testowego.

\section{Wyniki}

W obecnej wersji nie raportuje się wyników pomiarów, tabel statystycznych ani parametrów wyznaczonych z nieudokumentowanych danych. To kierunek badawczy, a nie zatwierdzony protokół badania ani zgłoszony artykuł.

\section{Dyskusja}

Interpretacja przyszłych wyników musi uwzględnić alternatywne mechanizmy i zależność wniosków od jakości danych, warunków zadania oraz ograniczeń sprzętu. Wartość potencjalnego rozwiązania zależy nie tylko od trafności modelu, lecz także od jego opóźnienia, stabilności i wpływu na zachowanie człowieka. Wyniki testów syntetycznych nie mogą być automatycznie uogólniane na populację uczestników.

\section{Ograniczenia badania}

Tobii Eye Tracker 5 jest urządzeniem ekranowym, nie okulometrem Quest 3; wymagane odrębne środowiska pomiaru. Granice nowości i trafności ekologicznej wymagają niezależnej weryfikacji. Żadna część szkicu nie dowodzi przeprowadzenia zbierania danych, uzyskania zgód etycznych ani gotowości manuskryptu do zgłoszenia.

\section{Wnioski}

Przedstawiono problem, kontrast i sposób jego planowanej oceny. Rozstrzygnięcie postawionego pytania wymaga przygotowania protokołu, zweryfikowanych źródeł oraz poprawnie udokumentowanych wyników. Bez tych materiałów nie można sformułować konkluzji empirycznych.

\begin{thebibliography}{9}
\bibitem{xrprior1} The integration of virtual reality and EEG: A step-by-step guideline (2026), Methods guideline. \url{https://pmc.ncbi.nlm.nih.gov/articles/PMC12808493/}.
\bibitem{xrprior2} Validity of Eye-Tracking Measures in Immersive Virtual Reality: Using Cognitive Theories of Learning (2026), Educational Psychology Review. \url{https://doi.org/10.1007/s10648-026-10186-9}.
\end{thebibliography}
\end{document}

